\section{Alineación humano-máquina y evaluación}

El expositor propone además la idea del ``agent judge'': en lugar de depender del etiquetado manual de todo el trace, se usa un agent optimizado por preferencias como judge, que califica el plan de cada interacción, y luego esas señales de preferencia se retroalimentan al proceso de generation.
\begin{table}[H]
\centering
\begin{tabular}{lp{6cm}}
\toprule
Métrica & Descripción \\
\midrule
Accuracy & Si los campos clave (destino, presupuesto, tiempo) son correctos y consistentes \\
Proactivity & Si formula suficientes preguntas dentro de un número limitado de turnos para completar el contexto \\
Credibility & Autoconsistencia y frecuencia de hallucination; el judge penaliza las respuestas que se contradicen a sí mismas \\
Diversity & Cantidad de planes en la conversación de varios turnos y cobertura de distintas estrategias \\
\bottomrule
\end{tabular}
\caption{Métricas de evaluación de agent judge}
\end{table}

\begin{knowledgebox}{Agent Judge+Human Alignment}
La salida del agent judge se somete a un soft alignment con proxies de human preference (como plan optimality y constraint compliance), y mediante DPO se retropropaga para que el generator converja más rápido hacia el comportamiento aprobado por el usuario.
\end{knowledgebox}

\subsection{Resumen de la sección}

El agent judge ofrece un bucle de retroalimentación de alta frecuencia: el reward de cada conversación se usa tanto para la actualización DPO del propio judge como para calibrar el planner, formando un ciclo cerrado de metrics + training.

